% !TEX root = ../tesi.tex
\chapter{Introduzione}
\label{chap:introduzione}

L'automazione della raccolta in serra tramite robot è uno dei temi più interessanti 
e attuali dell'agricoltura di precisione. In questo capitolo viene introdotto il contesto 
generale del lavoro, spiegando i motivi pratici ed etici per cui è utile automatizzare 
la raccolta dei pomodori. Successivamente viene descritto il problema principale 
affrontato nella tesi, ovvero la difficoltà di gestire le occlusioni all'interno dei 
grappoli densi, per poi presentare gli obiettivi principali del lavoro.


\section{Contesto e Motivazioni}
\label{sec:contesto_motivazioni}

Il lavoro svolto in questa tesi fa parte di un progetto di ricerca più ampio del 
Dipartimento di Informatica dell'Università degli Studi di Milano, coordinato dai 
professori Nicola Basilico e Matteo Luperto, con l'obiettivo di applicare la robotica 
e l'intelligenza artificiale all'agricoltura di precisione.

Negli ultimi anni la coltivazione in serra ha avuto una grande diffusione, poiché permette 
di controllare l'ambiente di crescita delle piante, risparmiare acqua, limitare l'uso di 
pesticidi e produrre durante tutto l'anno. Nonostante questi miglioramenti tecnologici, 
la gestione della serra richiede ancora molto lavoro manuale, soprattutto per il 
controllo delle piante e per la raccolta dei frutti \cite{gongal2015sensors}.

In particolare, la raccolta del pomodoro è una delle attività più faticose e costose. 
Gli operatori devono lavorare per molte ore all'interno delle serre, dove le temperature 
e i livelli di umidità sono spesso molto alti. A questo si aggiunge la crescente difficoltà 
nel trovare manodopera stagionale, oltre a una motivazione etica molto importante: 
l'automazione di questi compiti faticosi può contribuire a ridurre fenomeni di sfruttamento 
e caporalato, ancora troppo presenti nel settore agricolo. L'uso di robot per la raccolta 
permetterebbe inoltre di lavorare in modo continuo, ottimizzando i tempi di produzione.

In passato, i primi tentativi di automatizzare la raccolta si sono basati su macchine 
meccaniche come scuotitori o pettini. Questi sistemi funzionano a "ciclo aperto" 
(\textit{open-loop}): scuotono l'intera pianta senza distinguere quali frutti siano pronti 
e quali no, rischiando di rovinare la pianta e provocando molti sprechi. L'obiettivo della 
robotica moderna è invece realizzare sistemi a "ciclo chiuso" (\textit{closed-loop}), 
capaci di compiere una raccolta selettiva \cite{silwal2017design}: il robot deve raccogliere 
soltanto i frutti maturi al punto giusto per il mercato, lasciando intatti quelli acerbi 
in modo che possano maturare nei giorni successivi.

Per fare questo, la visione artificiale è la componente fondamentale del robot. L'ambiente 
della serra è complesso e variabile: la luce solare cambia continuamente durante il giorno, 
c'è polvere e la vegetazione è molto fitta. Prima di poter muovere il braccio meccanico per 
afferrare un pomodoro, il sistema deve essere in grado di "vedere" e capire la scena: 
deve individuare dove si trovano i frutti, capire quali sono maturi e quali acerbi, e 
verificare se siano facili da raggiungere, così da evitare urti accidentali con i rami o con la pianta.


\section{Definizione del Problema}
\label{sec:definizione_problema}

All'interno di una serra, il compito di raccogliere i pomodori si scontra con 
una difficoltà fondamentale: la struttura biologica della pianta. I pomodori 
non crescono isolati su rami sgombri, ma si sviluppano aggregati in 
grappoli densi e compatti, circondati da foglie, rami di sostegno e dal fusto principale.

In una scena così affollata, la sola identificazione visiva dei frutti tramite i modelli 
classici di Object Detection non è sufficiente per guidare un braccio robotico. Questi 
modelli racchiudono solitamente gli oggetti all'interno di riquadri rettangolari rigidi 
(\textit{bounding box}). Tuttavia, i pomodori hanno una sagoma tondeggiante o allungata: 
di conseguenza, i quattro angoli del rettangolo includono inevitabilmente una porzione 
consistente di spazio vuoto. Quando più pomodori crescono vicini o lungo linee oblique, 
i loro rettangoli di delimitazione finiscono per sovrapporsi, facendo credere al sistema 
che i frutti siano coperti tra loro anche quando in realtà sono completamente liberi.

Questa imprecisione geometrica rende molto complessa la gestione delle occlusioni reciproche, 
ovvero capire quale frutto si trovi davanti e quale dietro rispetto alla telecamera. Se il 
robot tentasse di raccogliere un pomodoro senza considerare chi lo copre frontalmente, 
la pinza di presa andrebbe a urtare i frutti antistanti, danneggiandoli. Inoltre, all'interno 
dello stesso grappolo sono spesso presenti pomodori ancora verdi: se un frutto acerbo si trova 
davanti a uno maturo, agisce come un ostacolo fisico che impedisce l'accesso al frutto rosso.

Infine, la raccolta di un grappolo non può essere pianificata in un colpo solo in modo statico. 
La rimozione fisica di un pomodoro modifica immediatamente la scena: togliere il primo frutto 
frontale libera spazio e permette di raggiungere i pomodori retrostanti che prima erano coperti. 
Per evitare collisioni e raccogliere i frutti nell'ordine corretto, è quindi necessario superare 
l'approssimazione dei rettangoli rigidi e sviluppare una strategia decisionale dinamica, capace 
di stimare i contatti reali e di aggiornare la sequenza di raccolta a ogni singolo prelievo.


\section{Obiettivi della Tesi}
\label{sec:obiettivi_tesi}

L'obiettivo principale di questo lavoro di tesi è progettare, implementare e 
validare una pipeline decisionale per stimare le occlusioni e pianificare la sequenza 
di raccolta dei pomodori in serra, superando i limiti delle tradizionali 
\textit{bounding box} attraverso una modellazione geometrica basata su ellissi orientate.

Nello specifico, il lavoro si articola sui seguenti obiettivi:
\begin{itemize}
    \item \textbf{Modellazione geometrica con Convex Hull ed ellissi}: elaborare le 
    maschere fornite dal modello di segmentazione YOLO11-x regolarizzandone i contorni 
    tramite l'involucro convesso e adattarvi un'ellisse orientata, eliminando lo spazio 
    vuoto negli angoli e stimando la reale inclinazione del frutto;
    \item \textbf{Screening e filtro di maturità}: valutare la raggiungibilità geometrica 
    dei pomodori (tramite area visibile, circolarità e centralità), eliminare le predizioni 
    duplicate e distinguere i frutti maturi da quelli acerbi tramite l'analisi del colore 
    nello spazio HSV, considerando i pomodori verdi come ostacoli da non urtare;
    \item \textbf{Task planning dinamico \textit{Pick \& Update}}: sviluppare un algoritmo 
    che simula il prelievo del frutto a ogni passo e aggiorna lo stato della scena, 
    premiando con un bonus di sblocco i pomodori retrostanti appena liberati;
    \item \textbf{Validazione sperimentale e confronto}: testare l'algoritmo su immagini 
    reali di serra con grappoli densi, confrontando le scelte del nostro modello sia con 
    una baseline a rettangoli rigidi sia con le decisioni prese da operatori umani.
\end{itemize}

